
%% bare_jrnl.tex
%% V1.3
%% 2007/01/11
%% by Michael Shell
%% see http://www.michaelshell.org/
%% for current contact information.
%%
%% This is a skeleton file demonstrating the use of IEEEtran.cls
%% (requires IEEEtran.cls version 1.7 or later) with an IEEE journal paper.
%%
%% Support sites:
%% http://www.michaelshell.org/tex/ieeetran/
%% http://www.ctan.org/tex-archive/macros/latex/contrib/IEEEtran/
%% and
%% http://www.ieee.org/



% *** Authors should verify (and, if needed, correct) their LaTeX system  ***
% *** with the testflow diagnostic prior to trusting their LaTeX platform ***
% *** with production work. IEEE's font choices can trigger bugs that do  ***
% *** not appear when using other class files.                            ***
% The testflow support page is at:
% http://www.michaelshell.org/tex/testflow/


%%*************************************************************************
%% Legal Notice:
%% This code is offered as-is without any warranty either expressed or
%% implied; without even the implied warranty of MERCHANTABILITY or
%% FITNESS FOR A PARTICULAR PURPOSE! 
%% User assumes all risk.
%% In no event shall IEEE or any contributor to this code be liable for
%% any damages or losses, including, but not limited to, incidental,
%% consequential, or any other damages, resulting from the use or misuse
%% of any information contained here.
%%
%% All comments are the opinions of their respective authors and are not
%% necessarily endorsed by the IEEE.
%%
%% This work is distributed under the LaTeX Project Public License (LPPL)
%% ( http://www.latex-project.org/ ) version 1.3, and may be freely used,
%% distributed and modified. A copy of the LPPL, version 1.3, is included
%% in the base LaTeX documentation of all distributions of LaTeX released
%% 2003/12/01 or later.
%% Retain all contribution notices and credits.
%% ** Modified files should be clearly indicated as such, including  **
%% ** renaming them and changing author support contact information. **
%%
%% File list of work: IEEEtran.cls, IEEEtran_HOWTO.pdf, bare_adv.tex,
%%                    bare_conf.tex, bare_jrnl.tex, bare_jrnl_compsoc.tex
%%*************************************************************************

% Note that the a4paper option is mainly intended so that authors in
% countries using A4 can easily print to A4 and see how their papers will
% look in print - the typesetting of the document will not typically be
% affected with changes in paper size (but the bottom and side margins will).
% Use the testflow package mentioned above to verify correct handling of
% both paper sizes by the user's LaTeX system.
%
% Also note that the "draftcls" or "draftclsnofoot", not "draft", option
% should be used if it is desired that the figures are to be displayed in
% draft mode.
%
\documentclass[draftclsnofoot,12pt,journal,onecolumn]{IEEEtran}
\renewcommand{\rmdefault}{phv} % Arial
\renewcommand{\sfdefault}{phv} % Arial
\setcounter{page}{3}

%
% If IEEEtran.cls has not been installed into the LaTeX system files,
% manually specify the path to it like:
% \documentclass[journal]{../sty/IEEEtran}





% Some very useful LaTeX packages include:
% (uncomment the ones you want to load)


% *** MISC UTILITY PACKAGES ***
%
%\usepackage{ifpdf}
% Heiko Oberdiek's ifpdf.sty is very useful if you need conditional
% compilation based on whether the output is pdf or dvi.
% usage:
% \ifpdf
%   % pdf code
% \else
%   % dvi code
% \fi
% The latest version of ifpdf.sty can be obtained from:
% http://www.ctan.org/tex-archive/macros/latex/contrib/oberdiek/
% Also, note that IEEEtran.cls V1.7 and later provides a builtin
% \ifCLASSINFOpdf conditional that works the same way.
% When switching from latex to pdflatex and vice-versa, the compiler may
% have to be run twice to clear warning/error messages.






% *** CITATION PACKAGES ***
%
\usepackage{cite}
% cite.sty was written by Donald Arseneau
% V1.6 and later of IEEEtran pre-defines the format of the cite.sty package
% \cite{} output to follow that of IEEE. Loading the cite package will
% result in citation numbers being automatically sorted and properly
% "compressed/ranged". e.g., [1], [9], [2], [7], [5], [6] without using
% cite.sty will become [1], [2], [5]--[7], [9] using cite.sty. cite.sty's
% \cite will automatically add leading space, if needed. Use cite.sty's
% noadjust option (cite.sty V3.8 and later) if you want to turn this off.
% cite.sty is already installed on most LaTeX systems. Be sure and use
% version 4.0 (2003-05-27) and later if using hyperref.sty. cite.sty does
% not currently provide for hyperlinked citations.
% The latest version can be obtained at:
% http://www.ctan.org/tex-archive/macros/latex/contrib/cite/
% The documentation is contained in the cite.sty file itself.



%\usepackage{hyperref}


% *** GRAPHICS RELATED PACKAGES ***
%
\ifCLASSINFOpdf
  % \usepackage[pdftex]{graphicx}
  % declare the path(s) where your graphic files are
  % \graphicspath{{../pdf/}{../jpeg/}}
  % and their extensions so you won't have to specify these with
  % every instance of \includegraphics
  % \DeclareGraphicsExtensions{.pdf,.jpeg,.png}
\else
  % or other class option (dvipsone, dvipdf, if not using dvips). graphicx
  % will default to the driver specified in the system graphics.cfg if no
  % driver is specified.
   \usepackage[dvips]{graphicx}
  % declare the path(s) where your graphic files are
  % \graphicspath{{../eps/}}
  % and their extensions so you won't have to specify these with
  % every instance of \includegraphics
  % \DeclareGraphicsExtensions{.eps}
\fi
% graphicx was written by David Carlisle and Sebastian Rahtz. It is
% required if you want graphics, photos, etc. graphicx.sty is already
% installed on most LaTeX systems. The latest version and documentation can
% be obtained at: 
% http://www.ctan.org/tex-archive/macros/latex/required/graphics/
% Another good source of documentation is "Using Imported Graphics in
% LaTeX2e" by Keith Reckdahl which can be found as epslatex.ps or
% epslatex.pdf at: http://www.ctan.org/tex-archive/info/
%
% latex, and pdflatex in dvi mode, support graphics in encapsulated
% postscript (.eps) format. pdflatex in pdf mode supports graphics
% in .pdf, .jpeg, .png and .mps (metapost) formats. Users should ensure
% that all non-photo figures use a vector format (.eps, .pdf, .mps) and
% not a bitmapped formats (.jpeg, .png). IEEE frowns on bitmapped formats
% which can result in "jaggedy"/blurry rendering of lines and letters as
% well as large increases in file sizes.
%
% You can find documentation about the pdfTeX application at:
% http://www.tug.org/applications/pdftex





% *** MATH PACKAGES ***
%
\usepackage[cmex10]{amsmath}
\usepackage{amsfonts}
% A popular package from the American Mathematical Society that provides
% many useful and powerful commands for dealing with mathematics. If using
% it, be sure to load this package with the cmex10 option to ensure that
% only type 1 fonts will utilized at all point sizes. Without this option,
% it is possible that some math symbols, particularly those within
% footnotes, will be rendered in bitmap form which will result in a
% document that can not be IEEE Xplore compliant!
%
% Also, note that the amsmath package sets \interdisplaylinepenalty to 10000
% thus preventing page breaks from occurring within multiline equations. Use:
%\interdisplaylinepenalty=2500
% after loading amsmath to restore such page breaks as IEEEtran.cls normally
% does. amsmath.sty is already installed on most LaTeX systems. The latest
% version and documentation can be obtained at:
% http://www.ctan.org/tex-archive/macros/latex/required/amslatex/math/





% *** SPECIALIZED LIST PACKAGES ***
%
%\usepackage{algorithmic}
% algorithmic.sty was written by Peter Williams and Rogerio Brito.
% This package provides an algorithmic environment fo describing algorithms.
% You can use the algorithmic environment in-text or within a figure
% environment to provide for a floating algorithm. Do NOT use the algorithm
% floating environment provided by algorithm.sty (by the same authors) or
% algorithm2e.sty (by Christophe Fiorio) as IEEE does not use dedicated
% algorithm float types and packages that provide these will not provide
% correct IEEE style captions. The latest version and documentation of
% algorithmic.sty can be obtained at:
% http://www.ctan.org/tex-archive/macros/latex/contrib/algorithms/
% There is also a support site at:
% http://algorithms.berlios.de/index.html
% Also of interest may be the (relatively newer and more customizable)
% algorithmicx.sty package by Szasz Janos:
% http://www.ctan.org/tex-archive/macros/latex/contrib/algorithmicx/




% *** ALIGNMENT PACKAGES ***
%
%\usepackage{array}
% Frank Mittelbach's and David Carlisle's array.sty patches and improves
% the standard LaTeX2e array and tabular environments to provide better
% appearance and additional user controls. As the default LaTeX2e table
% generation code is lacking to the point of almost being broken with
% respect to the quality of the end results, all users are strongly
% advised to use an enhanced (at the very least that provided by array.sty)
% set of table tools. array.sty is already installed on most systems. The
% latest version and documentation can be obtained at:
% http://www.ctan.org/tex-archive/macros/latex/required/tools/


%\usepackage{mdwmath}
%\usepackage{mdwtab}
% Also highly recommended is Mark Wooding's extremely powerful MDW tools,
% especially mdwmath.sty and mdwtab.sty which are used to format equations
% and tables, respectively. The MDWtools set is already installed on most
% LaTeX systems. The lastest version and documentation is available at:
% http://www.ctan.org/tex-archive/macros/latex/contrib/mdwtools/


% IEEEtran contains the IEEEeqnarray family of commands that can be used to
% generate multiline equations as well as matrices, tables, etc., of high
% quality.


%\usepackage{eqparbox}
% Also of notable interest is Scott Pakin's eqparbox package for creating
% (automatically sized) equal width boxes - aka "natural width parboxes".
% Available at:
% http://www.ctan.org/tex-archive/macros/latex/contrib/eqparbox/





% *** SUBFIGURE PACKAGES ***
\usepackage[tight,footnotesize]{subfigure}
% subfigure.sty was written by Steven Douglas Cochran. This package makes it
% easy to put subfigures in your figures. e.g., "Figure 1a and 1b". For IEEE
% work, it is a good idea to load it with the tight package option to reduce
% the amount of white space around the subfigures. subfigure.sty is already
% installed on most LaTeX systems. The latest version and documentation can
% be obtained at:
% http://www.ctan.org/tex-archive/obsolete/macros/latex/contrib/subfigure/
% subfigure.sty has been superceeded by subfig.sty.



%\usepackage[caption=false]{caption}
%\usepackage[font=footnotesize,caption=false]{subfig}
% subfig.sty, also written by Steven Douglas Cochran, is the modern
% replacement for subfigure.sty. However, subfig.sty requires and
% automatically loads Axel Sommerfeldt's caption.sty which will override
% IEEEtran.cls handling of captions and this will result in nonIEEE style
% figure/table captions. To prevent this problem, be sure and preload
% caption.sty with its "caption=false" package option. This is will preserve
% IEEEtran.cls handing of captions. Version 1.3 (2005/06/28) and later 
% (recommended due to many improvements over 1.2) of subfig.sty supports
% the caption=false option directly:
%\usepackage[caption=false,font=footnotesize]{subfig}
%
% The latest version and documentation can be obtained at:
% http://www.ctan.org/tex-archive/macros/latex/contrib/subfig/
% The latest version and documentation of caption.sty can be obtained at:
% http://www.ctan.org/tex-archive/macros/latex/contrib/caption/




% *** FLOAT PACKAGES ***
%
%\usepackage{fixltx2e}
% fixltx2e, the successor to the earlier fix2col.sty, was written by
% Frank Mittelbach and David Carlisle. This package corrects a few problems
% in the LaTeX2e kernel, the most notable of which is that in current
% LaTeX2e releases, the ordering of single and double column floats is not
% guaranteed to be preserved. Thus, an unpatched LaTeX2e can allow a
% single column figure to be placed prior to an earlier double column
% figure. The latest version and documentation can be found at:
% http://www.ctan.org/tex-archive/macros/latex/base/



%\usepackage{stfloats}
% stfloats.sty was written by Sigitas Tolusis. This package gives LaTeX2e
% the ability to do double column floats at the bottom of the page as well
% as the top. (e.g., "\begin{figure*}[!b]" is not normally possible in
% LaTeX2e). It also provides a command:
%\fnbelowfloat
% to enable the placement of footnotes below bottom floats (the standard
% LaTeX2e kernel puts them above bottom floats). This is an invasive package
% which rewrites many portions of the LaTeX2e float routines. It may not work
% with other packages that modify the LaTeX2e float routines. The latest
% version and documentation can be obtained at:
% http://www.ctan.org/tex-archive/macros/latex/contrib/sttools/
% Documentation is contained in the stfloats.sty comments as well as in the
% presfull.pdf file. Do not use the stfloats baselinefloat ability as IEEE
% does not allow \baselineskip to stretch. Authors submitting work to the
% IEEE should note that IEEE rarely uses double column equations and
% that authors should try to avoid such use. Do not be tempted to use the
% cuted.sty or midfloat.sty packages (also by Sigitas Tolusis) as IEEE does
% not format its papers in such ways.


%\ifCLASSOPTIONcaptionsoff
%  \usepackage[nomarkers]{endfloat}
% \let\MYoriglatexcaption\caption
% \renewcommand{\caption}[2][\relax]{\MYoriglatexcaption[#2]{#2}}
%\fi
% endfloat.sty was written by James Darrell McCauley and Jeff Goldberg.
% This package may be useful when used in conjunction with IEEEtran.cls'
% captionsoff option. Some IEEE journals/societies require that submissions
% have lists of figures/tables at the end of the paper and that
% figures/tables without any captions are placed on a page by themselves at
% the end of the document. If needed, the draftcls IEEEtran class option or
% \CLASSINPUTbaselinestretch interface can be used to increase the line
% spacing as well. Be sure and use the nomarkers option of endfloat to
% prevent endfloat from "marking" where the figures would have been placed
% in the text. The two hack lines of code above are a slight modification of
% that suggested by in the endfloat docs (section 8.3.1) to ensure that
% the full captions always appear in the list of figures/tables - even if
% the user used the short optional argument of \caption[]{}.
% IEEE papers do not typically make use of \caption[]'s optional argument,
% so this should not be an issue. A similar trick can be used to disable
% captions of packages such as subfig.sty that lack options to turn off
% the subcaptions:
% For subfig.sty:
% \let\MYorigsubfloat\subfloat
% \renewcommand{\subfloat}[2][\relax]{\MYorigsubfloat[]{#2}}
% For subfigure.sty:
% \let\MYorigsubfigure\subfigure
% \renewcommand{\subfigure}[2][\relax]{\MYorigsubfigure[]{#2}}
% However, the above trick will not work if both optional arguments of
% the \subfloat/subfig command are used. Furthermore, there needs to be a
% description of each subfigure *somewhere* and endfloat does not add
% subfigure captions to its list of figures. Thus, the best approach is to
% avoid the use of subfigure captions (many IEEE journals avoid them anyway)
% and instead reference/explain all the subfigures within the main caption.
% The latest version of endfloat.sty and its documentation can obtained at:
% http://www.ctan.org/tex-archive/macros/latex/contrib/endfloat/
%
% The IEEEtran \ifCLASSOPTIONcaptionsoff conditional can also be used
% later in the document, say, to conditionally put the References on a 
% page by themselves.





% *** PDF, URL AND HYPERLINK PACKAGES ***
%
\usepackage{url}
% url.sty was written by Donald Arseneau. It provides better support for
% handling and breaking URLs. url.sty is already installed on most LaTeX
% systems. The latest version can be obtained at:
% http://www.ctan.org/tex-archive/macros/latex/contrib/misc/
% Read the url.sty source comments for usage information. Basically,
% \url{my_url_here}.





% *** Do not adjust lengths that control margins, column widths, etc. ***
% *** Do not use packages that alter fonts (such as pslatex).         ***
% There should be no need to do such things with IEEEtran.cls V1.6 and later.
% (Unless specifically asked to do so by the journal or conference you plan
% to submit to, of course. )


% correct bad hyphenation here
\hyphenation{op-tical net-works semi-conduc-tor}


\begin{document}
%
% paper title
% can use linebreaks \\ within to get better formatting as desired
\title{Principles and results \\ (C4EU 5.6.1: Evaluation of Results and Best Practices -a )}
%
%
% author names and IEEE memberships
% note positions of commas and nonbreaking spaces ( ~ ) LaTeX will not break
% a structure at a ~ so this keeps an author's name from being broken across
% two lines.
% use \thanks{} to gain access to the first footnote area
% a separate \thanks must be used for each paragraph as LaTeX2e's \thanks
% was not built to handle multiple paragraphs
%

\author{
	Fernando~Gros, %\IEEEmembership{Member,~IEEE,}
	Alejandro~Andreu, %\IEEEmembership{Member,~IEEE,}
	Jorge~Beltran, %\IEEEmembership{Member,~IEEE,}
	Ignacio~Justel, %\IEEEmembership{Member,~IEEE,}
    and~Jaume~Barcelo% \IEEEmembership{Life~Fellow,~IEEE}% <-this % stops a space
\thanks{
Affiliation
}
}


% note the % following the last \IEEEmembership and also \thanks - 
% these prevent an unwanted space from occurring between the last author name
% and the end of the author line. i.e., if you had this:
% 
% \author{....lastname \thanks{...} \thanks{...} }
%                     ^------------^------------^----Do not want these spaces!
%
% a space would be appended to the last name and could cause every name on that
% line to be shifted left slightly. This is one of those "LaTeX things". For
% instance, "\textbf{A} \textbf{B}" will typeset as "A B" not "AB". To get
% "AB" then you have to do: "\textbf{A}\textbf{B}"
% \thanks is no different in this regard, so shield the last } of each \thanks
% that ends a line with a % and do not let a space in before the next \thanks.
% Spaces after \IEEEmembership other than the last one are OK (and needed) as
% you are supposed to have spaces between the names. For what it is worth,
% this is a minor point as most people would not even notice if the said evil
% space somehow managed to creep in.



% The paper headers
\markboth{C4EU 5.1.2: Report on Selection of Opportunities and Projects -b}%
{C4EU 5.1.2: Report on Selection of Opportunities and Projects -b}
% The only time the second header will appear is for the odd numbered pages
% after the title page when using the twoside option.
% 
% *** Note that you probably will NOT want to include the author's ***
% *** name in the headers of peer review papers.                   ***
% You can use \ifCLASSOPTIONpeerreview for conditional compilation here if
% you desire.




% If you want to put a publisher's ID mark on the page you can do it like
% this:
%\IEEEpubid{0000--0000/00\$00.00~\copyright~2007 IEEE}
% Remember, if you use this you must call \IEEEpubidadjcol in the second
% column for its text to clear the IEEEpubid mark.



% use for special paper notices
%\IEEEspecialpapernotice{(Invited Paper)}




% make the title area
\maketitle


\begin{abstract}
%\boldmath
This report specifies the driving principles of bottom up broadband initiatives.
These principles are developed to describe the best practices in terms of networking and communications, taking into account related reference models such as the open source community.
The report also includes a review of the project indicators and which is the current state of the project.
\end{abstract}
% IEEEtran.cls defaults to using nonbold math in the Abstract.
% This preserves the distinction between vectors and scalars. However,
% if the journal you are submitting to favors bold math in the abstract,
% then you can use LaTeX's standard command \boldmath at the very start
% of the abstract to achieve this. Many IEEE journals frown on math
% in the abstract anyway.

% Note that keywords are not normally used for peerreview papers.
\begin{IEEEkeywords}
 Bottom-up-Broadband (BuB), principles, results
%Slotted Aloha, game theory, contention control, media access control.
\end{IEEEkeywords}

\clearpage

\tableofcontents

\clearpage

\listoffigures

\listoftables

\clearpage




% For peer review papers, you can put extra information on the cover
% page as needed:
% \ifCLASSOPTIONpeerreview
% \begin{center} \bfseries EDICS Category: 3-BBND \end{center}
% \fi
%
% For peerreview papers, this IEEEtran command inserts a page break and
% creates the second title. It will be ignored for other modes.
\IEEEpeerreviewmaketitle



\section{Introduction}
% The very first letter is a 2 line initial drop letter followed
% by the rest of the first word in caps.
% 
% form to use if the first word consists of a single letter:
% \IEEEPARstart{A}{demo} file is ....
% 
% form to use if you need the single drop letter followed by
% normal text (unknown if ever used by IEEE):
% \IEEEPARstart{A}{}demo file is ....
% 
% Some journals put the first two words in caps:
% \IEEEPARstart{T}{his demo} file is ....
% 
% Here we have the typical use of a "T" for an initial drop letter
% and "HIS" in caps to complete the first word.
\IEEEPARstart{T}{his} report summarizes the principles that drive our work, the identified best practices and the results obtained so far.
As we promote collaboration, this document has been written in a collaborative way. 
The details as well as the instructions to collaborate are provided in Section \ref{sec:about}.
The core BuB principles are explained in Section \ref{sec:principles}.
As there are several parallelisms between BuB and open source communities, we devote Section \ref{sec:open-source} to explore them.
Wireless community networks have established a number of networking principles that foster collaboration, which are reviewed in Section \ref{sec:open-network}.

The communication and discussion of results are critical for long term sustainability in BuB initiatives and, for this reason, we emphasize the importance of documentation and transparency in Section \ref{sec:documentation} and some guidelines to prepare workshops in \ref{sec:workshops}.
We assess which is the current status of the project by means of the quantitative indicators examined in Section \ref{sec:indicators}.
Finally, Section \ref{sec:conclusion} concludes this report.

\section{About this document}
\label{sec:about}

This report has been produced using open source tools such as {\LaTeX} \cite{lamport1994ldp} and \emph{git} \cite{chacon2009pg}.
{\LaTeX} is widely used in academia to prepare print-class documents.
It automatically takes care of numbering, cross-referencing, tables of contents, bibliography, etc.
\emph{Git} is a high performance distributed revision control which is used in many open source projects, such as the linux kernel.
Git makes it easy and safe to collaborate as each contributor works on his own personal copy.
Good contributions can be easily shared with others, and it is always possible to revert to a previous version.

Our git repository is publicly available in \emph{github}:

https://github.com/jbarcelo/C4EU-deliverables

Anyone who is familiar with {\LaTeX} and \emph{github} can contribute to this document.
The firs step is to make a copy (a \emph{fork} in \emph{github} jargon).
The contributor can work in this copy and make changes to improve the document.
After that, it is necessary to request that these changes are merged into the original copy of the document (a \emph{pull request} in github jargon).

If you see anything that can be improved, feel free to contribute.
This document is alive in the sense that it will keep evolving as long as contributors make changes and improve it.

The system automatically keeps track of all the contributors and their contributions.
It is possible to see who is contributing more actively and which are the exact changes made by each contributor.
And everything is public on the web.

\section{BuB Principles}
\label{sec:principles}

Bottom up broadband has its starting point in an unsatisfied need by the users.
This users can be either individuals, communities, institutions, companies or other organizations.
The users get involved in the planning, funding and deployment of the solution.

Having a working solution is not the only goal in BuB initiatives.
Even though approaches such as fast prototyping are encouraged in BuB, BuB might not always be the fastest or more economical way.

In BuB there is also an interest in learning, in sharing and in collaborating.
BuB is about educating the users and empowering them.
For this reason, there is an emphasis on workshops, participative mailing lists and accessible documentation.

It is also important to be inclusive and understand that everyone can contribute in one way or another.
It is by no means restricted to experts or gurus.
Everyone has some strengths and weaknesses, and by means of collaboration it is possible to combine the strengths of each of us to achieve more ambitious goals.

As the African proverb says ``If you want to go fast, go alone. If you want to go far, go together.''

\subsection{Address issues that matter. Start with the people}

It is critical to focus the available resources to try to solve problems that matter.
We are interested in initiatives that will have an impact on people's lives, and therefore we need the people to help us to identify those initiatives.
As Tim Brown explains it in \cite{brown2009ted}, we want to ``start with humans''.
Integrate technology and economics, and start with what humans need.
Consequently, the first step has to be to understand culture and context, and working with the communities in the identification of relevant challenges.

Another principle emphasized in \cite{brown2009ted} is ``learning by making'', or ``building to think''.
Fast prototyping makes it possible to advance quicker, and to identify the strengths and weaknesses of the initiatives early on.
This same idea is sometime expressed as working for ``quick wins''.
It boils down to trying to achieve early results as a first step before attempting more ambitious goals.

A third idea pointed out in \cite{brown2009ted} is ``from consumption to participation''.
If we talk about government, this means stop complaining and fixing it.
This is the message that is conveyed by Jennifer Pahlka, founder of Code for America, in her TED talk ``Coding for a better government'' \cite{pahlka2012ted}.


\section{A parallelism with open source communities}
\label{sec:open-source}

The spirit of BuB is very similar to the open source initiative. Before diving into similarities, it is worth saying that the only significant difference is the user attitude. In a bottom up broadband context the user must provide some resources to become part of the resulting network, while in open source community both developers and users benefit from external contributions. Of course, both approaches are completely free.

Transparency is the most evident similarity with this kind of community. For an end user to connect to the network some specifications must be handled thus making this information public and transparent is a key aspect of the BuB approach. The more information, the more users. 

From transparency comes derived works. That is, individuals and organizations can adapt these projects to their needs.

As well as the open source definition states, any discrimination or restrictions take place in the BuB philosophy.




\section{Open Network}
\label{sec:open-network}

Open network is a network that allows free connection for anyone, i.e. no restrictions to prevent someone else connecting to it and therefore promotes the growth of the network.


The open network concept itself, is limited to the deployment and operation of the network regardless the content and services running on it. Therefore open network ensures that nobody can prohibit the entrance to new people and essential services network (IPs, DNS, etc.) works properly. An important aspect of the open network is that it allows connection to many segments of the population excluded (village with low population) since its deployment model reduces costs compared to the currently dominant models.


The guifi.net foundation is a clear example of open network, also with the participation of all users has created the largest open network of the world. It is an open network because all network settings data are published, allowing any person can see how the network is built, to improve, maintain and extend.


A misconception that people have about the open network is that it is not secure, because the link layer is not encrypted, allowing any user to analyze the data passing through your node, but in any network (internet or open network) security depends on the user (application layer), therefore the user is responsible for implementing security measures according to their needs. For example, protocols such as SSL allow secure channel connection between two points in the network and that all information that travels over the network is encrypted.


\section{Transparency and documentation}
\label{sec:documentation}

To achieve the goals set at the beginning of the project, it is paramount to have the necessary documentation to support the achievements and progress of the same. Not only for the fact of having it, but also to give help and support to potential new contributions from other members that look interested in the project. One of the qualities of Commons for Europe, is that it is completely and totally open, and anyone interested can contribute in a big or small way, depending on their knowledge or involvement. What is sought with this idea is to achieve that C4EU far from being a project of a team gets to be a project done by all.

One of the main objectives which we desire to achieve is the fact of get C4EU expanded and taken place in the whole European Union. To accomplish this objective, it is essential, therefore, the fact of being transparent in all our actions and knowledge and try to give as much help as possible to help new integrators to achieve the project objectives. Whether through direct contributions or by providing documentation in which we rely to carry it out. For more information please visit the official web site of the project or take a look out code at GitHub.com.


\section{Open Workshops Preparation}
\label{sec:workshops}

So far, there have been two workshops and both have been organized by people from the UPF. At this section, the details on how to prepare an open workshop will be provided.


As any other kind of meeting, some steps must be followed in order to organize a workshop:

\begin{enumerate}
	\item \textbf{Decide the need to make a workshop}
	\item \textbf{Let the partners know:} Once, you have decided that a workshop will take place it is recommended to inform the partners 					about this fact.
	\item \textbf{Agree on a date:} At this point, you have to ask your partners about their availability and take a poll in order to 							choose the best day for everyone.
	\item \textbf{Book a room:} Once the date and time are fixed, you may reserve a room for the workshop.
	\item \textbf{Inform about the requirements:} The information on the documents and/or other kind of things that have to be prepared in 					order to participate in the workshop should be provided at this point. You may also give a deadline for the documents to be 							uploaded and explain the different roles that any participant can take.  
	\item \textbf{Remote participation:} It is recommended to provide a way to attend to the workshop remotely, since may be some partners 					unable to attend physically.  
	\item \textbf{Attendance sheet:} Prepare an attendance sheet providing a summary of all the important information and any other data 						that have to be highlighted. All the participants should sign that sheet, indicating the role they will take. 
\end{enumerate}
 

\section{Project indicators}
\label{sec:indicators}

Some quantitative success indicators have been proposed.
Table \ref{tab:indicators} summarize the indicators that have been consider and the rest of the section offers further details on each of them.

\begin{table}[!t]
%% increase table row spacing, adjust to taste
\renewcommand{\arraystretch}{1.3}
% if using array.sty, it might be a good idea to tweak the value of
%\extrarowheight as needed to properly center the text within the cells
\caption{Success indicators}
\label{tab:indicators}
\centering
\begin{tabular}{|c|c|c|c|}
\hline
Indicator & Low & Medium & High \\
Pilot Proposals & 0 - 10 & 11 - 15 & +15 \\
Executed sensor pilots & 1 & 2 & +3 \\
Number of executed WiFi/SuperWiFi pilots& 1 & 2 & +3 \\
Number of executed fiber pilots& 1 & 2 & +3 \\
Number of executed hybrid pilots& 0 & 1 & 2 \\
BuB final users & 0-500 & 500-5000 & +5000 \\
BuB raised fundss & 0-1000 & 1000-100,000& +100,000 \\
Cities involved & 0-2 & 3-5 & +6 \\
\hline
\hline
\end{tabular}
\end{table}


\subsection{Pilot proposals 0-10/11-15/+15}

The leaders of the pilots work package will issue two calls for pilots
in 2012 and 2013. All the partners of the project and external BuB
champions will respond to this call with pilot proposals. These
proposals must contain some basic information regarding the project,
and clearly indicate which are the differentiating characteristics
that make them more attractive and suitable for execution. Among all
the proposed pilots, the partnership will choose those that are more
interesting following criteria of potential impact, geographical
distribution, technological diversity, and commitment of the champions
that are backing the proposal. The resources of the C4U project will
be devoted to monitoring, supporting and documenting the selected
pilots.

So far we have received 14 pilot proposals.

\subsection{Number of executed sensor pilots 1/2/3}
The most promising sensor pilots of all the the received proposals
will be executed. These pilots will receive full support from the
partners involved in the BuB pilots work package and will be carefully
documented to replicate the success models and disseminate the lessons
learned.

At this point, there is one sensor pilot being executed.

\subsection{Number of executed superwifi/wifi pilots 1/2/3}
The most promising superwifi and wifi pilots of all the the received
proposals will be executed. These pilots will receive full support
from the partners involved in the BuB pilots work package and will be
carefully documented to replicate the success models and disseminate
the lessons learned.

At this point, there are two WiFi pilots in execution.
One deals with public WiFi ofering while the other addresses mobile WiFi in combination with mesh technology.

\subsection{Number of executed fiber pilots 1/2/3}
The most promising fiber pilots of all the the received proposals will
be executed. These pilots will receive full support from the partners
involved in the BuB pilots work package and will be carefully
documented to replicate the success models and disseminate the lessons
learned.

There is one fiber pilot being executed at this moment.

\subsection{ Number of executed hybrid pilots 0/1/2}
Some of the pilots will combine more of one of the technologies of
interest in the C4EU project. This hybrid pilots represent a major
challenge as they will require the collaboration of different working
groups with heterogeneous background and expertise to successfully
integrate disparate technologies. Note that a hybrid project involving
two different technologies will be counted thrice. As an example, a
hybrid wifi-sensor pilot counts as a wifi pilot, a sensor pilot and a
hybrid pilot.

There are preliminar talks to launch an hybrid pilot which studies the impact of a fiber connection in an urban wireless network.

\subsection{BuB final users 0-500/500-5000/+5000}
One of the goals of this project is to have an impact on a large
number of citizens and institutions. To this end, we will estimate the
aggregate number of users of the pilots deployed in the project. As an
example, if a BuB project provides Internet connectivity to a cultural
center, we will estimate how many people (including staff and
visitors) have benefited from the pilot.

At this point, none of the pilots has been completed and therefore this metric is difficult to estimate.

\subsection{BuB raised funds 0-1000eur/1000-100,000eur/+100,000eur}
The BuB principle requires involvement and commitment from the part of
those that will benefit from the pilot. This commitment takes the form
of time and effort devoted to the initiative (which are difficult to
quantify) and financial support for the execution of the pilot. This
indicator reflects the amount of BuB money that is raised from sources
different than the project. These should cover the major part of the
expenses of the pilots using commercial technologies (wifi, fiber).

Some of the pilots have raised over 18,000 Euro in crowdfunding.
Nevertheless, the total amount of money will be difficult to estimate until the pilots are completed.
Figure \ref{fig:crowdfunded} shows an example of a crowdfunded initiative in the goteo.org crowdfunding platform.

\begin{figure}[htbp]
  \centering
  \includegraphics[width = 11cm]{figures/crowdfunded.eps}
  \caption{An example of a crowdfunded initiative in goteo.org (in Spanish)}
  \label{fig:crowdfunded}
\end{figure}

\subsection{Cities  involved in the pilots 0-2/3-5/+6}
One of the goals of the project is to geographically spread success
models across Europe, and to establish successful collaborations among
champions that share the same goals in different locations. This goal
will be fulfilled if the projects are executed in several different
cities.

So far we have pilots in Gurb (Vic), Barcelona, London, the cities of the province of Rome.
Some of the pilots are mobile or portable and can easily spread to a large number of cities.

\section{Conclusion}
\label{sec:conclusion}
In this report we reviewed the driving principles in BuB initiatives and how they develop in best practices.
An analysis of the quantitative project indicators shows that the project is in a good position to achieve the figures defined as medium and good performance for all indicators.




% if have a single appendix:
%\appendix[Proof of the Zonklar Equations]
% or
%\appendix  % for no appendix heading
% do not use \section anymore after \appendix, only \section*
% is possibly needed

% use appendices with more than one appendix
% then use \section to start each appendix
% you must declare a \section before using any
% \subsection or using \label (\appendices by itself
% starts a section numbered zero.)
%


%\appendices
%\section{Proof of the First Zonklar Equation}
%Appendix one text goes here.

% you can choose not to have a title for an appendix
% if you want by leaving the argument blank
%\section{}
%Appendix two text goes here.


% use section* for acknowledgement
\section*{Acknowledgment}

This work has been partially funded by the European Commission (grant CIP-ICT PSP-2011-5).
The views expressed in this technical report are solely those of the authors and do not represent the views of the European Commission.


% Can use something like this to put references on a page
% by themselves when using endfloat and the captionsoff option.
%\ifCLASSOPTIONcaptionsoff
%  \newpage
%\fi



% trigger a \newpage just before the given reference
% number - used to balance the columns on the last page
% adjust value as needed - may need to be readjusted if
% the document is modified later
%\IEEEtriggeratref{8}
% The "triggered" command can be changed if desired:
%\IEEEtriggercmd{\enlargethispage{-5in}}

% references section

% can use a bibliography generated by BibTeX as a .bbl file
% BibTeX documentation can be easily obtained at:
% http://www.ctan.org/tex-archive/biblio/bibtex/contrib/doc/
% The IEEEtran BibTeX style support page is at:
% http://www.michaelshell.org/tex/ieeetran/bibtex/
\bibliographystyle{IEEEtran}
% argument is your BibTeX string definitions and bibliography database(s)
\bibliography{IEEEabrv,my_bib}
%
% <OR> manually copy in the resultant .bbl file
% set second argument of \begin to the number of references
% (used to reserve space for the reference number labels box)
%\begin{thebibliography}{1}

%\bibitem{IEEEhowto:kopka}
%H.~Kopka and P.~W. Daly, \emph{A Guide to \LaTeX}, 3rd~ed.\hskip 1em plus
%  0.5em minus 0.4em\relax Harlow, England: Addison-Wesley, 1999.

%\end{thebibliography}

% biography section
% 
% If you have an EPS/PDF photo (graphicx package needed) extra braces are
% needed around the contents of the optional argument to biography to prevent
% the LaTeX parser from getting confused when it sees the complicated
% \includegraphics command within an optional argument. (You could create
% your own custom macro containing the \includegraphics command to make things
% simpler here.)
%\begin{biography}[{\includegraphics[width=1in,height=1.25in,clip,keepaspectratio]{mshell}}]{Michael Shell}
% or if you just want to reserve a space for a photo:

%\begin{IEEEbiography}{Michael Shell}
%Biography text here.
%\end{IEEEbiography}

% if you will not have a photo at all:
%\begin{IEEEbiographynophoto}{John Doe}
%Biography text here.
%\end{IEEEbiographynophoto}

% insert where needed to balance the two columns on the last page with
% biographies
%\newpage

%\begin{IEEEbiographynophoto}{Jane Doe}
%Biography text here.
%\end{IEEEbiographynophoto}

% You can push biographies down or up by placing
% a \vfill before or after them. The appropriate
% use of \vfill depends on what kind of text is
% on the last page and whether or not the columns
% are being equalized.

%\vfill

% Can be used to pull up biographies so that the bottom of the last one
% is flush with the other column.
%\enlargethispage{-5in}



% that's all folks
\end{document}


